\documentclass[10pt,a4paper]{article}
\usepackage[T1]{fontenc}
\usepackage[margin=2cm]{geometry}
\usepackage{isabelle,isabellesym}
\usepackage{amsmath,amssymb}
\usepackage{pdfsetup}

\urlstyle{rm}
\isabellestyle{it}

\begin{document}

\title{Perfect Difference Sets of Positive Integers}
\author{Carles Mar\'in}
\maketitle

\begin{abstract}
A set $A$ of positive integers is a perfect difference set if every nonzero integer is $a-b$ for exactly one
pair $(a,b)$ of distinct elements of $A$. Erd\H{o}s remarked that such sets are easy to construct~\cite{Er80},
and Lev gave a greedy construction~\cite{Lev04}: at each step add $z$ and $z+d$, where $d$ is the least positive
integer that is not yet a difference. We formalize the existence of a perfect difference set, with an explicit
choice of the new element, and the bound $a_n\le 70\,n^3$ for the greedy construction that takes the least
admissible new element, where $a_n$ is the larger element of the representation of $n$. Admissibility is phrased
through sums (the Sidon property), so that the forbidden values of the new element are at most $8m^3$ functions of
triples of the $m$ elements already chosen. Both results are stated as in the Formal Conjectures
project~\cite{FC} (Erd\H{o}s problem 1194).
\end{abstract}

\paragraph{Acknowledgement.} The formalization and the text of this document were prepared with the
assistance of generative AI systems, and checked by the author.

\tableofcontents

\parindent 0pt\parskip 0.5ex

\input{session}

\bibliographystyle{abbrv}
\bibliography{root}

\end{document}
